\section{C1 UI: a full-parent remote complement}
\label{sec:c1-ui-remote}
The source has actual ordinary Alice angle and infinitesimal Bob
angle. Apply the exact native stabilizer (AS5) as in (NI7) and take
$\eta$ a positive bounded unit, $\zeta=g^2$ infinitesimal and $k$ a
bounded unit. The reflected orientation need not satisfy $k\le1$.
Its $E$ remains a positive unit by the exact source entry in (NI7).
Normalize by the positive pure-source $d$.
Equations (NC2)--(NC4) give
\[
 C/d,\ A/d,\quad g^2m/d,\quad g^2\beta/d,\quad
 \alpha/(dg^2)\quad\hbox{strong},\qquad D/d\ge0.       \tag{CR1}
\]
For example $P_d,P_\alpha,Q,r+\eta$ are units, whereas the
positive openings $M,t,r,1-r$ are bounded and may be infinitesimal.
The $-1/2$ in $\beta/d$ disappears after multiplication by $g^2$.
In this particular source subdomain $\beta=0$ is impossible; $D=0$
is allowed throughout. Choose positive real constants
$\mu,\nu,c,a,b$ below the five quantities in (CR1), respectively
in the order $g^2m/d,g^2\beta/d,C/d,A/d,\alpha/(dg^2)$.
Whenever needed we cap these coefficients at
\[
 m=\mu/g^2,\quad\beta=\nu/g^2,\quad C=c,\quad A=a,
 \quad\alpha=bg^2.                                  \tag{CR2}
\]
Every unused packet, every positive excess, and the complete $D$
packet is retained as a square. The signed core, divided by two, is
\[
 K=J^2+D_u(vY-Z)^2+D_v(X+uZ)^2-D_uD_vZ^2,
 \quad J=uW+uvX-Y-vZ.                              \tag{CR3}
\]
Only the core and the explicitly selected caps need a bounded Gram.
This observation keeps all upper coefficient scales out of the proof.

\begin{theorem}[C1 UI remote tests]
For every source satisfying (CR1), the complete original C1 parent
is tensor SOS at every independent testing pair except possibly
$u$ a bounded nonzero unit and $v$ infinitesimal. The conclusion
includes literal projective endpoints, both signs, all states,
and $D=0$.
\end{theorem}
\begin{proof}
All constants in residue forms below are real residues. Roots may be
unbounded, but all coefficients in their exact Gram are bounded.
First put $Q=W+uY$, $B=X+uZ$; then
\[
 M=vQ-B,\quad uW-Y=uQ-D_uY,\quad
 J=uQ-D_uY+uvB-vD_uZ.                               \tag{CR4}
\]
These identities account for every state normal.

If $u$ is a bounded unit and $v$ projectively infinite, put $w=1/v$
and homogenize by $w^2$. Use roots
$(E_Q,E_B,Y,Z)=(Q/g,B/g,Y,Z)$ and cap only $m,\beta$.
The exact homogenized core is
\[
 [w(ugE_Q-D_uY)+ugE_B-D_uZ]^2
 +D_u(Y-wZ)^2+D_wg^2E_B^2-D_uD_wZ^2.
\]
Its residue is $D_u(Y^2+u^2Z^2)$. Including the two caps gives
\[
 4\mu E_Q^2+4\nu E_B^2+2D_u(Y^2+u^2Z^2)>0.         \tag{CR5}
\]

If $u$ is infinitesimal and $v$ projectively infinite, the preceding
four-root frame loses the $Z$ pivot. Keep $w=1/v$ and instead use
$(E_Q,E_B,Y,U,T)$, where $U=uZ$ and $T=gwZ$. Put
$A_0=w(ugE_Q-D_uY)+ugE_B$. The exact homogenized core is
\[
 K_w=A_0^2+D_uY^2+D_wg^2E_B^2+D_uU^2
       -2D_ug(wE_Q+E_B)U+2wD_uuYU.
\]
All its coefficients are bounded, with residue $Y^2+U^2$.
Cap $m,\beta,\alpha$; their amplitudes after normalization are
$E_Q-wE_B$, $E_B$, and $gY-T$, with the positive factors
$D_w,D_u$ in the last two packets. The full selected residue is
\[
 4\mu E_Q^2+4\nu E_B^2+2Y^2+2U^2+4bT^2>0.       \tag{CR5a}
\]
This includes literal $u=0$ or $w=0$; no division by either was used.

If $u$ is infinitesimal and $v$ a bounded unit, use
$E_Q=Q/g,E_B=B/g,Y,U=uZ,Z_g=gZ$, and cap $m,\beta,\alpha$.
Substitute $W=gE_Q-uY$, $X=gE_B-uZ$ in (CR3).
All terms involving $Z$ collect into $uZ$; the resulting exact
core is a polynomial in $(u,v,g;E_Q,E_B,Y,U)$.
Its residue is $D_vY^2+v^2U^2$. The full selected residue is
\[
 4\mu(vE_Q-E_B)^2+4\nu D_vE_B^2
       +2D_vY^2+2v^2U^2+4bZ_g^2>0.                 \tag{CR6}
\]
The $\alpha$ amplitude is $gvY-Z_g$, so no inverse of $u$ was used
in constructing this Gram.

If both $u,v$ are infinitesimal, set
\[
 U=uQ,\quad V=vQ/g,\quad E_B=B/g,\quad Z_g=gZ,
 \quad S=uvZ,
\]
and include $Y$ as the sixth root. The exact core is
\begin{align*}
K={}&(U-D_uY+uvgE_B)^2+D_uv^2Y^2+D_vg^2E_B^2+D_uS^2\\
 &+2D_uuYS-2D_uuVZ_g-2D_uvgE_BS.                 \tag{CR7}
\end{align*}
In particular the apparently problematic term $UvZ$ is exactly
$uVZ_g$, with no mixed inverse. Cap all five weights in (CR2).
The entire selected residue is
\[
 4\mu(V-E_B)^2+4\nu E_B^2+(4c+2)(U-Y)^2
          +4aY^2+4bZ_g^2+2S^2>0.                  \tag{CR8}
\]

It remains to cover $u$ projectively infinite. Put $p=1/u$,
homogenize by $p^2$, and use $D_p=1+p^2$,
$Q=pW+Y$, $B=pX+Z$. Then
\[
 M=vQ-B,\quad r=W-pY,\quad
 J=D_p(W+vX)-p(Q+vB).                              \tag{CR9}
\]
For $v$ a bounded unit, use $(Q/g,B/g,W,X)$ and cap $m,\beta,C$.
The core after $Y=gE_Q-pW$, $Z=gE_B-pX$ is a bounded polynomial
with residue $(W+vX)^2$. The selected residue is
\[
 4\mu(vE_Q-E_B)^2+4\nu D_vE_B^2
           +4cD_vW^2+2(W+vX)^2>0.                 \tag{CR10}
\]

For infinitesimal $v$, use $(W,R,Y,V,E_B)$ with
$R=vX$, $V=vQ/g$, $E_B=B/g$. Cap $m,\beta,C,A$, retaining
$\alpha$ entirely. The exact core is
\begin{align*}
K_p={}&(W+D_pR-pY-pvgE_B)^2+D_vg^2E_B^2\\
 &+D_p\{v^2Y^2-2vgYE_B+2pYR-(vgE_B-pR)^2\}.
                                                               \tag{CR11}
\end{align*}
Every coefficient is bounded. The selected residue is
\[
 4\mu(V-E_B)^2+4\nu E_B^2+4cW^2+4aY^2
                    +2(W+R)^2>0.                 \tag{CR12}
\]

Finally, when both testing angles are infinite, also put $w=1/v$
and homogenize by $p^2w^2$. Use $(Q/g,B/g,W,X)$ and cap
$m,\beta,C$. Substitution $Y=gE_Q-pW$, $Z=gE_B-pX$ in the
homogenized core is polynomial and has residue $X^2$.
The selected residue is
\[
 4\mu E_Q^2+4\nu E_B^2+4cW^2+2X^2>0.               \tag{CR13}
\]
Each displayed positive form is strict on the complete root space.
The bounded positive-residue Gram lemma therefore supplies tensor
SOS. Restoring the retained packets proves the original parent.
The bounded unit/unit testing chart is already covered by the
inherited noncommuting-residue theorem. This leaves exactly the
ordinary/small testing strip stated in the theorem. The exact gauges
$(u,W,X)\mapsto(-u,-W,-X)$ and
$(v,W,Y)\mapsto(-v,-W,-Y)$ preserve every event square and the core.
All homogeneous formulas allow $p=0$ or $w=0$ literally.
Real and imaginary state components are treated separately.
\end{proof}

The checker verifies the exact core reductions and their residues;
it does not infer positivity from samples. This theorem narrows the
actual UI subcase of C1.5. It does not close C1.5, whose actual UH
source subcase remains, and it does not yet close UI's ordinary/small
strip. The leading-contact calculations stored alongside this proof
remain candidate material until their complete normal lift is printed.
